
\subsection{情景结果与机制讨论}
以下结论仅在当前参数组与 2022 公报归一化口径下成立，不构成独立外部验证。在该口径下，江豚末值和长江鲟末值分别为 1.133 和 1.140。若仅提高通道阻隔而不加污染，江豚和长江鲟分别降至 0.901 和 0.856，显示附件 1 对应的通道约束本身就会压低恢复收益。进一步将长江鲟人工放流项置为 $R_S=0$ 时，长江鲟末值降至 0.519，而江豚仅变为 1.145，显示人工放流主要抬升鲟类，对江豚不是主解释项。污染情景中，江豚和长江鲟分别降至 0.623 和 0.706，提示污染会在高阻隔基础上继续侵蚀基线情景中的恢复收益。这里不把 $M$ 和 $V$ 放进结果表，是因为它们是传导和竞争状态，不是本问的直接观测目标。
\begin{table}[htbp]
\centering
\caption{问题二的结果口径、相对变化与对照说明}
\begin{tabularx}{0.95\textwidth}{>{\raggedright\arraybackslash}p{2.2cm} >{\raggedright\arraybackslash}p{3.6cm} >{\raggedright\arraybackslash}p{2.8cm} >{\raggedright\arraybackslash}p{1.7cm} X}
\toprule
项目 & 数值/情景输出 & 相对基线变化 & 是否外部验证 & 解释用途 \\
\midrule
2022 公报锚点 & 江豚 1249 头、长江鲟 438 尾 & -- & 是 & 作为问题二的初值与口径锚定 \\
基线情景末值 & 江豚 1.133，长江鲟 1.140 & -- & 否 & 作为禁渔与食源修复下的基线投影 \\
$R_S=0$ 对照 & 江豚 1.145，长江鲟 0.519 & 江豚 +1.0\%，长江鲟 -54.5\% & 否 & 拆分自然恢复与人工放流补偿效应 \\
高阻隔对照 & 江豚 0.901，长江鲟 0.856 & 江豚 -20.5\%，长江鲟 -24.9\% & 否 & 仅提高通道阻隔，分离生境约束效应 \\
污染情景末值 & 江豚 0.623，长江鲟 0.706 & 江豚 -45.0\%，长江鲟 -38.0\% & 否 & 比较污染叠加阻隔对基线收益的侵蚀 \\
\bottomrule
\end{tabularx}
\end{table}
这组对照表明，放流主要抬升鲟类，污染主要压低两类珍稀物种，因此问题二应收口为“情景比较与机制提示”，而不是独立外部验证。

\begin{figure}[htbp]
\centering
\includegraphics[width=0.88\textwidth]{fig02_q2_rare_species.png}
\caption{问题二和问题五共用的稀有种与入侵模块结果}
\label{fig:q2}
\end{figure}
